\documentclass[a4paper]{article}
\usepackage[utf8]{inputenc}
\usepackage[english]{babel}
\usepackage{amsmath} % per ambienti tipo cases
\usepackage{mathtools}
\usepackage{siunitx}
\usepackage{graphicx} % per includere figure
%\usepackage{subfigure}
\usepackage{booktabs} % per le tabelle
\usepackage{caption}
\usepackage{fancyhdr}
\usepackage{hyperref}
\usepackage[section]{placeins}
\usepackage{microtype}
\usepackage{caption}
\usepackage{subcaption}
\captionsetup[subfigure]{labelfont=rm}
\usepackage{verbatim} %multiline comments
%\usepackage[backend=biber, style=numeric, safeinputenc, sorting=none]{biblatex}
%\addbibresource{source.bib}	% uncomment for bibliography



%opening
\title{}
\author{}

\pagestyle{fancy}
\lhead{Musical Acoustics}
\chead{Homework 2}
\rhead{10743504, 10751919}
\newcommand{\Rarrow}{\mbox{\Large$\Rightarrow$}}

\begin{document}

\begin{titlepage}	
	\newcommand{\HRule}{\rule{\linewidth}{0.5mm}} % Defines a new command for horizontal lines, change thickness here
	
	\center % Centre everything on the page
	
	%------------------------------------------------
	%	Headings
	%------------------------------------------------
	
	\includegraphics[width=.4\textwidth]{Logo_Politecnico_Milano.png}\\[0.4cm]
	\textsc{\LARGE}\\[0.3cm] % Main heading such as the name of your university/college
	
	\textsc{\large MSc. Music and Acoustic Engineering}\\[1cm] % Minor heading such as course title
	
	\textsc{\Large Musical Acoustics - A.Y. 2020/2021}\\[0.5cm] % Major heading such as course name
	
	%------------------------------------------------
	%	Title
	%------------------------------------------------
	
	\HRule\\[0.4cm]
	
	{\huge\bfseries Homework 3 - Sound Radiation from Plates}\\[0.4cm] % Title of your document
	
	\HRule\\[1.5cm]
	
	
	
	{\large\textit{Authors' IDs:}}\\
	10743504, 10751919, % Your name
	%\\ \textsc{Gruppo 11}
	
	%------------------------------------------------
	%	Date
	%------------------------------------------------
	
	\vfill\vfill\vfill % Position the date 3/4 down the remaining page
	
	{\large\today} % Date, change the \today to a set date if you want to be precise
	
	%------------------------------------------------
	%	Logo
	%------------------------------------------------
	
	\vfill\vfill
	%\includegraphics[width=0.2\textwidth]{Politecnico_di_Milano.eps}\\[1cm] % Include a department/university logo - this will require the graphicx package
	
	%----------------------------------------------------------------------------------------
	
	\vfill % Push the date up 1/4 of the remaining page
	
	
\end{titlepage}

\section{Adjusting the radiation cutoff by choosing the plate thickness}

Assuming an infinite plate, we can study the sound radiation by imposing that the air's velocity field and the plate velocity are equal on the plate's surface. This yields the following expression for the component of the acoustic wave vector normal to the surface:
$$ k_z = \omega \sqrt{\frac{1}{c^2} - \frac{1}{v_p^2}} $$
where $c$ is the sound velocity in air and $v_p$ is the velocity of the bending waves in the plate. Since bending waves in a plate are dispersive ($v_p \propto \omega^{\frac{1}{2}}$), the argument of the square root increases with frequency, which means there'll be a cutoff frequency below which $k_z$ becomes imaginary. In this case the solution for the acoustic field takes the form of a vanishing wave in the z direction, which doesn't carry any energy away from the plate's surface. Therefore, below the cutoff frequency there will not be any sound radiation.

In particular, the cutoff frequency is the one for which $v_p = c$. The velocity of the bending waves in a thin plate is:
$$ v_p(\omega) = \sqrt{\frac{\omega hc_L}{\sqrt{12}}} $$
where $c_L$ is the corresponding velocity of the quasi-longitudinal waves. For a wooden plate the longitudinal velocity depends on the direction of the wave; indicating with $x$ and $y$ the longitudinal and radial axes of the material respectively: 
\begin{align*}
	c_x = \sqrt{E_x /\rho(1 - \nu_{xy}\nu_{yx})} \qquad& c_y = \sqrt{E_y /\rho(1 - \nu_{yx}\nu_{xy})}
\end{align*}




\end{document}